\chapter{Radical Prostatectomy}
\section*{Overview}
Radical prostatectomy removes the prostate and seminal vesicles, with reconnection of the bladder to the urethra, to treat selected localized prostate cancer.
\section*{Common Instruments and Equipment}
The exact setup varies by technique and setting. Common items may include:
\begin{itemize}[nosep,leftmargin=*]
  \item open, laparoscopic, or robotic pelvic surgery instruments
  \item urinary catheter and anastomosis sutures
  \item specimen retrieval bag for minimally invasive surgery
\end{itemize}

\section*{Why It Is Done}
It is considered for men with cancer confined or largely confined to the prostate after discussion of surveillance, radiation, and other options.
\section*{How It Is Performed}
The operation may be open, laparoscopic, or robotic. Pelvic lymph nodes may be removed for staging, and nerve-sparing may be attempted when oncologically appropriate.
\section*{Contraindications}
No single contraindication applies to every patient; the treating team assesses whether the procedure is suitable for the individual and the reason for treatment.

\section*{Precautions}
The care team reviews relevant health conditions, medicines, allergies, and procedure-specific risks, and explains any changes needed for the individual.

\section*{Risks and Recovery}
Urinary incontinence, erectile dysfunction, bleeding, infection, urethral narrowing, lymphocele, and cancer recurrence are possible. PSA monitoring and rehabilitation are required.
\section*{References}
\begin{enumerate}\item National Cancer Institute. Prostate Cancer Surgery. 2025.\item Current Medical Diagnosis and Treatment. McGraw-Hill; 2025.\end{enumerate}
\clearpage
